\documentclass[10pt,a4paper]{amsart}
\usepackage{amsmath,amssymb,amsthm,graphicx}
\usepackage{verbatim}
\newenvironment{reference}{\comment}{\endcomment}
%\newenvironment{reference}{}{}
\newenvironment{slogan}{\comment}{\endcomment}
\newenvironment{history}{\comment}{\endcomment}

% For commutative diagrams we use Xy-pic
\usepackage[all]{xy}

% We use 2cell for 2-commutative diagrams.
\xyoption{2cell}
% Theorem environments.
%
\theoremstyle{plain}
\newtheorem{theorem}[subsection]{Theorem}
\newtheorem{proposition}[subsection]{Proposition}
\newtheorem{lemma}[subsection]{Lemma}

\theoremstyle{definition}
\newtheorem{definition}[subsection]{Definition}
\newtheorem{example}[subsection]{Example}
\newtheorem{exercise}[subsection]{Exercise}
\newtheorem{situation}[subsection]{Situation}

\theoremstyle{remark}
\newtheorem{remark}[subsection]{Remark}
\newtheorem{remarks}[subsection]{Remarks}

\numberwithin{equation}{subsection}

% Macros
%
\def\lim{\mathop{\mathrm{lim}}\nolimits}
\def\colim{\mathop{\mathrm{colim}}\nolimits}
\def\Spec{\mathop{\mathrm{Spec}}}
\def\Hom{\mathop{\mathrm{Hom}}\nolimits}
\def\Ext{\mathop{\mathrm{Ext}}\nolimits}
\def\SheafHom{\mathop{\mathcal{H}\!\mathit{om}}\nolimits}
\def\SheafExt{\mathop{\mathcal{E}\!\mathit{xt}}\nolimits}
\def\Sch{\mathit{Sch}}
\def\Mor{\mathop{\mathrm{Mor}}\nolimits}
\def\Ob{\mathop{\mathrm{Ob}}\nolimits}
\def\Sh{\mathop{\mathit{Sh}}\nolimits}
\def\NL{\mathop{N\!L}\nolimits}
\def\CH{\mathop{\mathrm{CH}}\nolimits}
\def\proetale{{pro\text{-}\acute{e}tale}}
\def\etale{{\acute{e}tale}}
\def\QCoh{\mathit{QCoh}}
\def\Ker{\mathop{\mathrm{Ker}}}
\def\Im{\mathop{\mathrm{Im}}}
\def\Coker{\mathop{\mathrm{Coker}}}
\def\Coim{\mathop{\mathrm{Coim}}}

% Boxtimes
%
\DeclareMathSymbol{\boxtimes}{\mathbin}{AMSa}{"02}

%
% Macros for moduli stacks/spaces
%
\def\QCohstack{\mathcal{QC}\!\mathit{oh}}
\def\Cohstack{\mathcal{C}\!\mathit{oh}}
\def\Spacesstack{\mathcal{S}\!\mathit{paces}}
\def\Quotfunctor{\mathrm{Quot}}
\def\Hilbfunctor{\mathrm{Hilb}}
\def\Curvesstack{\mathcal{C}\!\mathit{urves}}
\def\Polarizedstack{\mathcal{P}\!\mathit{olarized}}
\def\Complexesstack{\mathcal{C}\!\mathit{omplexes}}
% \Pic is the operator that assigns to X its picard group, usage \Pic(X)
% \Picardstack_{X/B} denotes the Picard stack of X over B
% \Picardfunctor_{X/B} denotes the Picard functor of X over B
\def\Pic{\mathop{\mathrm{Pic}}\nolimits}
\def\Picardstack{\mathcal{P}\!\mathit{ic}}
\def\Picardfunctor{\mathrm{Pic}}
\def\Deformationcategory{\mathcal{D}\!\mathit{ef}}
\usepackage{cleveref}
\crefname{originallemma}{original lemma}{original lemmas}
\Crefname{originallemma}{Original lemma}{Original lemmas}
\makeatletter
\newlabel{algebra-lemma-exact-sequence-NL}{{00S2}{source}{}{sourcepointer-1}{}}
\newlabel{algebra-lemma-exact-sequence-NL@cref}{{[originallemma][0][]00S2}{[1][0][]source}{}{}{}}
\makeatother
\crefname{originalreference}{original reference}{original references}
\Crefname{originalreference}{Original reference}{Original references}
\makeatletter
\newlabel{lemma-trace-well-defined}{{0ADZ}{source}{}{sourcepointer-2}{}}
\newlabel{lemma-trace-well-defined@cref}{{[originalreference][0][]0ADZ}{[1][0][]source}{}{}{}}
\makeatother
\title{Trace maps on higher differential forms for a degree p purely inseparable ring extension}
\author{The Stacks Project Authors}
\date{stacks-project:resolve.tex:local-snapshot-20261012; GNU FDL 1.2 or later}
\begin{document}
\maketitle
\noindent\textit{Editorial note.} All delivered mathematical text is the author's own, copied from the pinned Stacks Project chapter file resolve.tex; no statement or proof step was rewritten, reordered or removed.\par\medskip
\noindent\textit{Editorial note.} The counted proof is the author's complete argument: the degree one composition defining the trace, the computation of the cokernel of the connecting map through the transitivity exact sequence, and the exterior power construction with the explicit formula.\par\medskip
\noindent\textit{Editorial note.} The label lemma-trace-higher names the counted statement, and the proof environment immediately following it in the source is delivered verbatim as its proof.\par\medskip
\noindent\textit{Editorial note.} The counted statement is carried by the author's own lemma environment in the chapter source.\par\medskip
\noindent\textit{Editorial note.} Pre-open gate: python operations/agent-authoring/tools/precard\_gate.py --source resolve.tex --label lemma-trace-higher --start 500 --end 564 --local-id 28101 returned RESULT=OK (L28101 free, occupancy FREE, public interval clear). Pre-assembly tag self-check: canonical tag resolved and not present in reference/public-items.jsonl.\par\medskip
\noindent\textit{Editorial note.} The statement and proof are copied from the pinned chapter file; the cited inputs are kept as external pointers with their canonical Stacks tags.\par\medskip
\noindent\textit{Allowed standard prerequisites (editorial).}\begin{itemize}
\item The trace map on differentials for the presentation by a p-th power relation. Same chapter, the lemma constructing the trace map on differential forms in the degree one case. Supplies the degree one trace map whose image lies in the module generated by the differential of the relation, which is what makes the wedge construction well defined.
\item The transitivity exact sequence for differentials of a ring extension. Algebra, the exact sequence relating the first homology of the cotangent complex of a relative extension, the extended differentials of the base, the differentials of the source, and the relative differentials. Identifies the cokernel of the connecting map with the tensor product of a computable module, which is the input to the wedge construction.
\item The relative differentials and the first homology in the presentation case. The computation of the relative differentials as free on the differential of the generator, and of the first homology of the cotangent complex as free on the relation. Makes the connecting map explicit and hence the cokernel computable.
\end{itemize}\medskip
\ifdefined\hypertarget\hypertarget{sourcepointer-1}{}\fi
\noindent\textit{Source reference.} Original lemma 00S2: sources/stacks-chow/tags.txt, source lines 1024--1024. Prerequisite pointer invoked at proof line(s) 538; it names a standard fact of a sibling chapter that the counted argument uses, and every step of the counted proof built on it is retained verbatim in the delivered body.\par\medskip
\ifdefined\hypertarget\hypertarget{sourcepointer-2}{}\fi
\noindent\textit{Source reference.} Original reference 0ADZ: sources/stacks-resolve/pinned-original-source.tex, source lines 400--400. Prerequisite pointer invoked at proof line(s) 532; it names a standard fact of the same chapter that the counted argument uses, and every step of the counted proof built on it is retained verbatim in the delivered body.\par\medskip
\begin{lemma}
\label{lemma-trace-higher}
Let $\mathbf{F}_p \subset \Lambda \subset R \subset S$ be ring extensions
and assume that $S$ is isomorphic to $R[x]/(x^p - a)$ for some $a \in R$.
Then there are canonical $R$-linear maps
$$
\text{Tr} :
\Omega^{t + 1}_{S/\Lambda}
\longrightarrow
\Omega_{R/\Lambda}^{t + 1}
$$
for $t \geq 0$ such that
$$
\eta_1 \wedge \ldots \wedge \eta_t \wedge x^i\text{d}x
\longmapsto
\left\{
\begin{matrix}
0 & \text{if} & 0 \leq i \leq p - 2, \\
\eta_1 \wedge \ldots \wedge \eta_t \wedge \text{d}a & \text{if} & i = p - 1
\end{matrix}
\right.
$$
for $\eta_i \in \Omega_{R/\Lambda}$ and such that $\text{Tr}$ annihilates the
image of
$S \otimes_R \Omega_{R/\Lambda}^{t + 1} \to \Omega_{S/\Lambda}^{t + 1}$.
\end{lemma}

\begin{proof}
For $t = 0$ we use the composition
$$
\Omega_{S/\Lambda} \to \Omega_{S/R} \to \Omega_R \to \Omega_{R/\Lambda}
$$
where the second map is Lemma \ref{lemma-trace-well-defined}.
There is an exact sequence
$$
H_1(L_{S/R}) \xrightarrow{\delta} \Omega_{R/\Lambda} \otimes_R S \to
\Omega_{S/\Lambda} \to \Omega_{S/R} \to 0
$$
(Algebra, Lemma \ref{algebra-lemma-exact-sequence-NL}).
The module $\Omega_{S/R}$ is free over $S$ with basis $\text{d}x$
and the module $H_1(L_{S/R})$ is free over $S$ with basis $x^p - a$
which $\delta$ maps to $-\text{d}a \otimes 1$ in
$\Omega_{R/\Lambda} \otimes_R S$. In particular, if we set
$$
M = \Coker(R \to \Omega_{R/\Lambda}, 1 \mapsto -\text{d}a)
$$
then we see that $\Coker(\delta) = M \otimes_R S$. We obtain a
canonical map
$$
\Omega^{t + 1}_{S/\Lambda} \to
\wedge_S^t(\Coker(\delta)) \otimes_S \Omega_{S/R} =
\wedge^t_R(M) \otimes_R \Omega_{S/R}
$$
Now, since the image of the map
$\text{Tr} : \Omega_{S/R} \to \Omega_{R/\Lambda}$
of Lemma \ref{lemma-trace-well-defined} is contained in $R\text{d}a$ we
see that wedging with an element in the image annihilates $\text{d}a$.
Hence there is a canonical map
$$
\wedge^t_R(M) \otimes_R \Omega_{S/R} \to \Omega_{R/\Lambda}^{t + 1}
$$
mapping
$\overline{\eta}_1 \wedge \ldots \wedge \overline{\eta}_t \wedge \omega$
to $\eta_1 \wedge \ldots \wedge \eta_t \wedge \text{Tr}(\omega)$.
\end{proof}

\end{document}
